% Appendix H: Perturbative reduction in higher-curvature Lagrangians (~2100 w).
%
% Supersedes the "Constraint promotion in higher-curvature PGT" first draft.
% Sources (read-only, no verbatim reuse):
%   docs/tex/perturbative_reduction.tex          (§§1–2: EFT problem, v6 algorithm)
%   docs/tex/perturbative_reduction_constraint_barrier.tex   (§§3–4: LPS, barrier)

\section{Perturbative reduction in higher-curvature Lagrangians}%
\label{PerturbativeReduction}

\TIDAL{} treats a correction $\varepsilon\,\mathcal{L}_1$ to a well-posed base
Lagrangian $\mathcal{L}_0$ as a perturbative shift to the base-theory dynamics.
The effective-field-theory background, the equations-of-motion reduction, and the
Lagrangian reduction used to obtain a Hamiltonian are described here.

\paragraph*{Higher-order terms as effective corrections}
For theories whose Lagrangian admits the split
$\mathcal{L} = \mathcal{L}_0 + \varepsilon\,\mathcal{L}_1$ with $\varepsilon$ a small
dimensionless coupling, $\varepsilon\,\mathcal{L}_1$ is treated throughout as a
perturbative correction within the regime $\varepsilon\ll1$.

Such higher-order operators are the low-energy residues of an unknown ultraviolet
completion in the effective-field-theory
interpretation~\cite{Simon:1990ic,Simon:1990jn,Donoghue:1994dn,burgess2003quantum}.
Grosse-Knetter~\cite{Grosse-Knetter:1993tae,ostrogradsky1850memoires,Woodard:2015zca}
proves that expressing the higher time derivatives in $\mathcal{L}_1$ by substituting
the base-theory equations of motion is equivalent to a canonical field redefinition
$\varphi^{\mathcal{I}} \to \varphi^{\mathcal{I}} + \varepsilon\,\delta\varphi^{\mathcal{I}}$
that preserves the Poisson bracket structure.
The original higher-derivative Lagrangian and the perturbatively reduced one are
therefore physically equivalent within the perturbative regime.
This equivalence requires the base-theory equation of motion for each reduced field to
take the form $\partial_t^2\varphi^{\mathcal{I}} = F^{\mathcal{I}}[\varphi,\partial_t\varphi,
\nabla\varphi]$, since the substitution replaces every occurrence of
$\partial_t^2\varphi^{\mathcal{I}}$ in $\mathcal{L}_1$ with the right-hand
side~\cite{Grosse-Knetter:1993tae}. Fields whose base-theory equation is purely
algebraic (in the sense of \cref{SymbolicComputation}) violate the requirement and the
canonical equivalence breaks down.

The alternative --- treating $\mathcal{L}$ as an exact theory rather than a truncated
expansion --- encounters a structural obstruction.
Ostrogradsky's theorem~\cite{ostrogradsky1850memoires,Pais:1950za,Woodard:2015zca} states
that any non-degenerate Lagrangian whose Euler--Lagrange equations are of differential
order greater than two has a Hamiltonian unbounded below. The construction promotes
each time-derivative order up to $\partial_t^{N-1}\varphi$ to an independent canonical
coordinate, so a Lagrangian of order $N$ has $N$ conjugate momenta. Woodard's
review~\cite{Woodard:2015zca} shows explicitly that the resulting Hamiltonian is linear
in all but one of these momenta (the exception is the momentum conjugate to the
highest-order coordinate, which enters quadratically). The energy is therefore unbounded
below in either sign of any of the linearly-appearing momenta, and the corresponding
mode is a ghost.%
\footnote{Apparent counterexamples in gravity --- the Gauss--Bonnet term in four
spacetime dimensions~\cite{Lovelock:1971yv} and the higher-derivative terms in the
Einstein--Hilbert Lagrangian~\cite{hobson2006general,wald1984general} --- are excluded:
the former is topological and contributes only a boundary term, and the latter reduce to
a total divergence after integration by parts, leaving the Euler--Lagrange equations of
order two in both cases.}
The theorem is non-perturbative; the ghost is present for any non-zero $\varepsilon$.
The ghost mode has characteristic frequency $\sim 1/\sqrt{\varepsilon}$, exactly the
scale at which $\varepsilon\,\mathcal{L}_1$ becomes comparable to $\mathcal{L}_0$ and
the perturbative expansion itself breaks down. The ghost is therefore an artefact of
extrapolating the truncated Lagrangian beyond its regime of validity, not a physical
degree of freedom, and a perturbative treatment that respects the validity bound
$\varepsilon \omega t \ll 1$ never encounters it.
We therefore treat $\mathcal{L}_0$ as the base theory and $\varepsilon\,\mathcal{L}_1$
perturbatively.

\paragraph*{Perturbative reduction of the equations of motion}
We solve the equations of motion by the iterative order-reduction method of Parker
and Simon~\cite{parker1993reduced}, generalised to the multi-field linearised-PDE
regime of \TIDAL{}.

Using $q = \{\varphi^{\mathcal{I}}\}$ to denote the collection of dynamical fields
and writing the full equation of motion as $E[q;\,\varepsilon] = 0$~\cite{parker1993reduced,
jaen1986reduction},
the base operator is $E_{\mathrm{base}}[q] \equiv E[q;\,\varepsilon{=}0]$ and the correction operator
is $E_{\mathrm{corr}}[q] \equiv \partial_\varepsilon E[q;\,\varepsilon]\big|_{\varepsilon=0}$
from $\mathcal{L}_1$, which may carry higher-order time and spatial derivatives.
Expanding $q = q^{(0)} + \varepsilon\,q^{(1)} + \cdots$ and collecting by order:
\begin{align}
  \mathcal{O}(\varepsilon^0) &:\quad
    E_{\mathrm{base}}\bigl[q^{(0)}\bigr] = 0, \notag \\[2pt]
  \mathcal{O}(\varepsilon^1) &:\quad
    E_{\mathrm{base}}\bigl[q^{(1)}\bigr] = -E_{\mathrm{corr}}\bigl[q^{(0)}\bigr].
  \label{EOMHierarchy}
\end{align}
At every order the left-hand-side operator is the ghost-free second-order base operator
$E_{\mathrm{base}}$; the higher-derivative contributions from $\mathcal{L}_1$ appear
only as sources evaluated on the previous-order solution.
We call the order-$\varepsilon^0$ solution \emph{Pass~0} and the
order-$\varepsilon^1$ correction \emph{Pass~1}.

Two alternative reduction paths were considered. The
Ja{\'e}n--Llosa--Molina (JLM) substitution~\cite{jaen1986reduction}, which substitutes the
order-zero Euler--Lagrange equation directly into the higher-derivative correction in
the Lagrangian, has no substitution target when a field equation is purely algebraic at
order zero, since no $\partial_t^n\varphi^{\mathcal{I}}$ is available to solve for. The
Ostrogradsky auxiliary-field reduction~\cite{ostrogradsky1850memoires,Woodard:2015zca}
introduces each time-derivative order as a separate dynamical variable
$Q_n^{\mathcal{I}} \equiv \partial_t^n\varphi^{\mathcal{I}}$ for $n = 0, 1, \ldots, N-1$,
with $N$ the order of the highest time derivative in $\mathcal{L}$, yielding a first-order
system. This is not a perturbative method: it reformulates the full
$\varepsilon$-finite theory rather than truncating it, and the Ostrogradsky ghost
discussed above is therefore part of the dynamics. Direct numerical evolution of the
auxiliary system samples the ghost's unbounded-below spectrum and produces runaway
trajectories that lie outside the validity of the truncated effective Lagrangian.

\paragraph*{Modal derivation and the Duhamel kernel}
On flat periodic domains, the Pass~0 system at wavenumber $k$ reduces to the per-mode
ODE $\partial_t\hat{y}^{(0)} = M(k)\,\hat{y}^{(0)}$ (cf.~\cref{ModalPerMode}), where
$\hat{y}^{(0)}$ stacks the Fourier amplitudes of every field and velocity slot.
For the perturbative-reduction path, the eigendecomposition $M(k) = V D V^{-1}$ is
used rather than the Pad\'e scaling-and-squaring path of \cref{ModalSolver}, since the
per-eigenmode expansion is required explicitly for the Duhamel integral
\cref{DuhamelIntegral}.
With $D = \mathrm{diag}(\lambda_1,\lambda_2,\ldots)$ collecting the eigenvalues of $M(k)$
and $V$ the right-eigenvector matrix, the initial data $\hat{y}^{(0)}(0)$ is projected
into the eigenbasis as $c = V^{-1}\hat{y}^{(0)}(0)$; the diagonal evolution operator
$e^{Dt}$ then acts on $c$ componentwise, and $V$ rotates back:
\begin{equation}
  \hat{y}^{(0)}(t) = V\,e^{Dt}\,c.
  \label{PassZeroSolution}
\end{equation}

Pass~1 satisfies the inhomogeneous system
$\partial_t\hat{y}^{(1)} = M(k)\,\hat{y}^{(1)} + C_{\mathrm{src}}\,\hat{y}^{(0)}(t)$
with $\hat{y}^{(1)}(0) = 0$, where $C_{\mathrm{src}}$ is the Fourier representation of
the correction operator $E_{\mathrm{corr}}$ that appears as the order-$\varepsilon^1$
source in \cref{EOMHierarchy}.
Substituting $\hat{y}^{(1)} = V\hat{z}^{(1)}$ and the Pass~0 solution
$\hat{y}^{(0)}(t) = V e^{Dt}c$ into this system and multiplying through by $V^{-1}$
diagonalises the homogeneous part ($V^{-1}M(k)V = D$) and rotates the source into the
eigenbasis ($V^{-1}C_{\mathrm{src}}V \equiv S$), giving
\begin{equation}
  \partial_t\hat{z}^{(1)}
  \;=\;
  D\,\hat{z}^{(1)}
  \;+\; S\,e^{Dt}\,c,
  \qquad \hat{z}^{(1)}(0) = 0.
  \label{EigenODESystem}
\end{equation}
Component-wise this is a set of decoupled scalar inhomogeneous ODEs,
\begin{equation}
  \partial_t\hat{z}^{(1)}_i
  \;=\;
  \lambda_i\,\hat{z}^{(1)}_i
  \;+\; \sum_j S_{ij}\,c_j\,e^{\lambda_j t},
  \qquad \hat{z}^{(1)}_i(0) = 0,
  \label{EigenODE}
\end{equation}
the homogeneous term carrying the eigenvalue $\lambda_i$ of mode $i$ and the source term
expressing the rotated coupling of every Pass~0 mode $j$ into the Pass~1 evolution of
mode $i$.
Multiplying through by $e^{-\lambda_i t}$ converts the left side to
$\tfrac{d}{dt}[e^{-\lambda_i t}\hat{z}^{(1)}_i]$; integrating both sides from $0$ to
$t$~\cite{hochbruck2010exponential,almohy2011action} then gives
\begin{equation}
  e^{-\lambda_i t}\hat{z}^{(1)}_i(t)
  \;=\;
  \sum_j S_{ij}\,c_j
  \int_0^t e^{(\lambda_j - \lambda_i)\tau}\,\mathrm{d}\tau,
  \label{DuhamelIntegral}
\end{equation}
in which the exponential integrand evaluates in closed form.
Restoring the factor $e^{\lambda_i t}$ gives
\begin{equation}
  \hat{z}^{(1)}_i(t)
  \;=\;
  \sum_j S_{ij}\,c_j\,G(\lambda_i,\lambda_j;\,t),
  \label{DuhamelSolution}
\end{equation}
where the Duhamel kernel $G$ is~\cite{hochbruck2010exponential,almohy2011action}
\begin{equation}
  G(\lambda,\mu;\,t)
  \;=\;
  \begin{cases}
    \dfrac{e^{\mu t} - e^{\lambda t}}{\mu - \lambda} & \mu \neq \lambda, \\[6pt]
    t\,e^{\lambda t} & \mu = \lambda,
  \end{cases}
  \label{DuhamelKernel}
\end{equation}
Restoring the original basis gives the closed-form Pass~1 state,
\begin{equation}
  \hat{y}^{(1)}(t)
  \;=\;
  V\,\hat{z}^{(1)}(t)
  \;=\;
  \sum_i \left[\,\sum_j S_{ij}\,c_j\,G(\lambda_i,\lambda_j;\,t)\right] V_{:i},
  \label{PassOneSolution}
\end{equation}
where $V_{:i}$ denotes the $i$-th right eigenvector of $M(k)$. Pass~1 is therefore
computed entirely from the Pass~0 eigensystem; no additional differential equation
need be solved.
Near the degenerate limit $|(\mu-\lambda)t| \to 0$, direct evaluation of
\cref{DuhamelKernel} suffers catastrophic cancellation; setting $z = (\mu-\lambda)t$,
\TIDAL{} uses the Taylor expansion~\cite{almohy2011action}
\begin{equation}
  G\bigl|_{|z|<10^{-5}}
  \;=\;
  t\,e^{\lambda t}\sum_{j=0}^{12}\frac{z^{j}}{(j+1)!},
  \label{DuhamelTaylor}
\end{equation}
the threshold and term count following~\cite{almohy2011action}.

The iterative approximation holds while the cumulative phase error satisfies
$\varepsilon\,\omega\,t_{\mathrm{end}} \ll 1$~\cite{parker1993reduced}, monitored
at run time by \TIDAL{}'s validity diagnostic.

\subsection{Perturbative reduction of the Hamiltonian}
Energy measurement requires a Hamiltonian, and the standard Legendre transform applies
once the Lagrangian contains at most first-order time derivatives.
Integration by parts shifts derivatives between factors but cannot reduce the total
time-derivative count of a quadratic-in-fields term. Quadratic terms whose total
time-derivative order exceeds two therefore remain irreducible to a Lagrangian with at
most first-order time derivatives on every factor.

To reduce the Lagrangian before the Legendre transform, \TIDAL{} performs the
\emph{Lagrangian perturbative substitution}: the order-$\varepsilon^0$ equations of
motion are substituted into the higher-derivative terms of $\varepsilon\,\mathcal{L}_1$
~\cite{jaen1986reduction,Cheyette1988,Grosse-Knetter:1993tae}.
The order-$\varepsilon^0$ equations are solved simultaneously for the second time
derivative of each field,
\begin{equation}
  \partial_t^2 q_i
  \;=\;
  F_i\!\bigl[q,\,\partial_t q,\,\nabla q\bigr]
  \;+\; \mathcal{O}(\varepsilon),
  \label{LPSSubstitution}
\end{equation}
where $F_i$ depends only on first-order time derivatives and spatial derivatives.
Higher spatial derivatives in $\varepsilon\,\mathcal{L}_1$ do not affect the Legendre
transform; only higher time derivatives require reduction.
Where the correction introduces terms with time derivatives of order greater than two,
the rule \cref{LPSSubstitution} is extended by differentiating with respect to $t$ as
many times as needed, substituting \cref{LPSSubstitution} for any $\partial_t^2 q_j$
that appears at each step.
The resulting substitution is applied to the order-$\varepsilon^1$ Lagrangian, and the
series is truncated at $\mathcal{O}(\varepsilon^2)$; the reduced Lagrangian carries at
most first-order time derivatives and is admissible for the standard Legendre transform
of \cref{EulerLagrange}.
The substitution is faithful exactly when the JLM regularity
condition~\cite{jaen1986reduction} holds: the order-zero equation of motion for each
reduced field must admit a $\partial_t^2\varphi^{\mathcal{I}}$ that can be isolated and
substituted; equivalently, the leading-order kinetic matrix must be invertible.
Algebraic equations of motion fail this condition, and the substitution then silently
drops a promoted degree of freedom.

\paragraph*{Constraint promotion and perturbative-reduction failure}
The case where the JLM condition fails arises physically as constraint promotion. In
the Einstein--Cartan limit $\varepsilon=0$, certain metric and torsion components
satisfy algebraic relations and carry no kinetic term; turning on $\varepsilon$
infinitesimally adds higher-derivative kinetic content from $\varepsilon\,\mathcal{L}_1$
and promotes these relations to dynamical equations. The constraint structure of the
theory therefore changes discontinuously across $\varepsilon = 0$, and an order-by-order
expansion in $\varepsilon$ is blind to this promotion: the new mode is absent at
$\varepsilon = 0$ and cannot be reached by a power-series expansion.

This phenomenon has been mapped in the Poincar\'e gauge theory literature: the
nonlinear Hamiltonian analysis of the quadratic-torsion family establishes the
constraint structure of these theories~\cite{barker2021hamiltonian}, and the same
physics appears as the ``strong-coupling problem'' of the PGT
literature~\cite{barker2023higgs}, in which ``the rank of the matrix of constraint
brackets changes on singular surfaces'' of phase space.

Three independent reduction approaches reach the same conclusion. The
JLM substitution~\cite{jaen1986reduction} cannot execute its rule
without a $\partial_t^2\varphi^{\mathcal{I}}$ to isolate. Grosse-Knetter's
canonical-equivalence result~\cite{Grosse-Knetter:1993tae} requires the same
substitution rule to define the field redefinition, and so applies under the same
prerequisite. Glavan~\cite{Glavan2018} compares the JLM method, EFT
field redefinitions, and the Cayuso--Lehner non-perturbative scheme side by side, and
shows that all three break down at constraint promotion for the same architectural
reason: the leading-order kinetic matrix must be invertible. This Lagrangian-level
substitution failure --- not the Ostrogradsky ghost discussed earlier, which the
perturbative reduction of \cref{EOMHierarchy} evades by $\varepsilon$-truncation in
every higher-derivative theory (including the well-behaved Klein--Gordon biharmonic
case validated in \cref{Validation}) --- is what distinguishes the constraint-promoted
PGT theories from cases where the perturbative reduction succeeds.

\paragraph*{Structural analogues}
Constraint promotion belongs to a wider class of discontinuous limits in gravity and
gauge theory~\cite{deRham2014,hinterbichler2012theoretical}. The van Dam--Veltman--Zakharov
discontinuity~\cite{vandam1970massive,zakharov1970linearized} shows that the massless
limit of a Fierz--Pauli graviton recovers only five of the six general-relativistic
polarisations: the sixth arises through a constraint promoted at non-zero mass.
The Boulware--Deser non-decoupling~\cite{boulware1972gravitation} extends this to the
non-linear theory by identifying a ghost-like sixth mode once the Fierz--Pauli tuning
is broken.
The Maxwell--Proca boundary illustrates the same pattern: at finite vector-boson mass
the algebraic relation $\rcD{_\mu}\FieldA{^\mu} = 0$ holds as a constraint, while at
zero mass it becomes a gauge symmetry; no perturbative expansion in the mass encodes
the transition. The Stückelberg mechanism~\cite{RueggRuizAltaba2003} smooths this
constraint structure across the limit through an auxiliary field.
A non-perturbative Hamiltonian analysis of the pure $R^2$ analogue has been carried
out~\cite{barker2024r2}; the natural direction for \TIDAL{} is to extend the
perturbative-reduction framework of this appendix to cover the constraint-promotion
case --- for example, by introducing Stückelberg-type auxiliary fields that smooth the
$\varepsilon \to 0$ limit --- and so to recover a perturbative Hamiltonian for the
quadratic-PGT $\tilde{R}^2$ theories.

\paragraph*{Failure detection}
For the Hamiltonian, \TIDAL{} attempts the Lagrangian perturbative substitution
symbolically at derivation time.
Failure is detected by checking whether the order-$\varepsilon^0$ momentum
kernel changes rank when $\varepsilon$ is set to a non-zero test value; a rank
jump signals constraint promotion and the software emits a diagnostic that
marks the theory as having no reduced Hamiltonian.
The three affected theories --- \texttt{graviton\_torsion},
\texttt{torsion\_gertsenshtein}, and \texttt{torsion\_gertsenshtein\_combined},
all carrying $\tilde{R}^2$ --- retain their unreduced Hamiltonians, and
their energy diagnostics are recorded separately rather than alongside
conversion results.

\paragraph*{Implications for the surveys}
The EOM and Hamiltonian reductions have different scopes, and the surveys of
\cref{Results} are structured accordingly.
The iterative EOM reduction of \cref{EOMHierarchy} succeeds for all surveyed
theories, including the $\tilde{R}^2$ quadratic-PGT theories: conversion
probabilities are computed from field trajectories evolved by Pass~0 and
Pass~1, and no Hamiltonian is required for that computation.
Constraint promotion applies only to the Hamiltonian side and only
for the three theories named above; their energy diagnostics are deferred to future
work. The EOM evolution itself is well-defined in these cases, so one option is to
evaluate the unreduced Hamiltonian on the Pass~0/Pass~1 trajectory and accept the
$\mathcal{O}(\varepsilon^2)$ energy error from the unreduced higher-derivative content;
whether the constraint-promoted dynamics produce a measurable energy signature beyond
that error is an empirical question for future investigation.
Theories without higher-curvature terms are unaffected and use the standard
linearised methodology of \cref{SymbolicComputation} throughout.
